% Part: computability
% Chapter: computability-theory
% Section: applications-fixed-point

\documentclass[../../../include/open-logic-section]{subfiles}

\begin{document}

\olfileid{cmp}{thy}{apf}
\olsection{کاربرد قضیهٔ نقطهٔ ثابت}

قضیهٔ نقطهٔ ثابت اساساً امکان می‌دهد توابع محاسبه‌پذیر جزئی را برحسب
نمایه‌هایشان تعریف کنیم. برای نمونه، می‌توان نمایه‌ای چون $e$ یافت
که به ازای هر $y$،
\[
\cfind{e}(y) = e + y.
\]
به‌عنوان نمونه‌ای دیگر، می‌توان از برهان قضیهٔ نقطهٔ ثابت برای طراحی
برنامه‌ای در جاوا یا ++C استفاده کرد که خودش را چاپ می‌کند.

به یاد آورید اگر به ازای هر $e$، $W_e$ را دامنهٔ $\cfind{e}$ بگیریم،
دنبالهٔ $W_0$، $W_1$، $W_2$،~\dots مجموعه‌های محاسبه‌پذیرِ برشمردنی
را برمی‌شمارد. برخی از این مجموعه‌ها محاسبه‌پذیرند. می‌توان پرسید آیا
الگوریتمی هست که مقدار $x$ را ورودی بگیرد و اگر $W_x$ محاسبه‌پذیر
بود، نمایه‌ای برای تابع مشخصهٔ آن بازگرداند. پاسخ «نه» است؛ چنین
الگوریتمی وجود ندارد:

\begin{thm}
هیچ تابع محاسبه‌پذیر جزئیِ $f$ با ویژگی زیر وجود ندارد: هرگاه $W_e$
محاسبه‌پذیر باشد، $f(e)$ تعریف شده باشد و $\cfind{f(e)}$ تابع مشخصهٔ
آن باشد.
\end{thm}

\begin{proof}
بگذارید $f$ هر تابع محاسبه‌پذیر جزئی باشد؛ $e$ای می‌سازیم که $W_e$
محاسبه‌پذیر است، اما $\cfind{f(e)}$ تابع مشخصهٔ آن نیست. با استفاده
از قضیهٔ نقطهٔ ثابت می‌توانیم نمایه‌ای چون $e$ بیابیم که
\[
\cfind{e}(y) \simeq
\begin{cases}
0 & \text{اگر $y=0$ و $\cfind{f(e)}(0) \fdefined = 0$} \\
\text{تعریف‌نشده} & \text{در غیر این صورت.}
\end{cases}
\]
یعنی $e$ با اعمال قضیهٔ نقطهٔ ثابت بر تابع زیر به دست می‌آید:
\[
g(x,y) \simeq
\begin{cases}
0 & \text{اگر $y=0$ و $\cfind{f(x)}(0) \fdefined = 0$} \\
\text{تعریف‌نشده} & \text{در غیر این صورت.}
\end{cases}
\]
به‌طور غیررسمی می‌توان محاسبه‌پذیری جزئیِ $g$ را چنین دید: الگوریتم
روی ورودی‌های $x$ و~$y$ نخست بررسی می‌کند آیا $y=0$. اگر چنین باشد،
الگوریتم $f(x)$ را محاسبه می‌کند و سپس با ماشین جهانی
$\cfind{f(x)}(0)$ را محاسبه می‌کند. اگر محاسبهٔ آخر متوقف شود و~$0$
بازگرداند، الگوریتم~$0$ را بازمی‌گرداند؛ وگرنه متوقف نمی‌شود.

اکنون توجه کنید اگر $f(e)$ تعریف‌نشده باشد، آنگاه $W_e=\emptyset$ و
خودِ شرط لازمِ تعریف‌شدن $f(e)$ نقض شده است. اگر $f(e)$ تعریف شده و
$\cfind{f(e)}(0)$ تعریف شده و برابر~$0$ باشد، $\cfind{e}(y)$ دقیقاً
وقتی تعریف شده است که $y=0$؛ پس $W_e=\{0\}$. اگر $f(e)$ تعریف شده
باشد اما $\cfind{f(e)}(0)$ تعریف‌نشده یا نابرابر با~$0$ باشد، آنگاه
$W_e=\emptyset$. در دو حالت اخیر $\cfind{f(e)}$ تابع مشخصهٔ~$W_e$
نیست، زیرا روی ورودی~$0$ پاسخ نادرست می‌دهد.
\end{proof}

\end{document}
